% Préambule commun des schémas TikZ — charte Liouaï
% Chaque schéma commence par % Préambule commun des schémas TikZ — charte Liouaï
% Chaque schéma commence par % Préambule commun des schémas TikZ — charte Liouaï
% Chaque schéma commence par \input{_preambule} puis \begin{document}.
\documentclass[tikz,border=6pt]{standalone}
\usepackage[T1]{fontenc}
\usepackage[utf8]{inputenc}
\usepackage[scaled=0.95]{helvet}
\renewcommand{\familydefault}{\sfdefault}
\usepackage{amsmath,amssymb}
\usepackage[european]{circuitikz}
\usetikzlibrary{arrows.meta,positioning,calc,angles,quotes,decorations.pathreplacing,backgrounds}

% Couleurs du logo Liouaï Digital Education
\definecolor{navy}{HTML}{243743}
\definecolor{vert}{HTML}{28B78D}
\definecolor{vertfonce}{HTML}{167A5C}
\definecolor{vertclair}{HTML}{EAF7F2}
\definecolor{gris}{HTML}{A1A9AE}
\definecolor{gristexte}{HTML}{5C6870}
\definecolor{fond}{HTML}{F3F5F6}
\definecolor{jaune}{HTML}{E5C800}
\definecolor{jauneclair}{HTML}{FFFBD6}

\tikzset{
  every picture/.style={line cap=round, line join=round, color=navy, text=navy},
  fleche/.style={-{Stealth[length=2.6mm,width=2mm]}, line width=1.1pt},
  vecteur/.style={fleche, line width=1.4pt},
  titre/.style={font=\bfseries\small, text=navy},
  note/.style={font=\footnotesize, text=gristexte},
  boite/.style={draw=navy, line width=0.9pt, rounded corners=5pt, fill=fond, align=center, inner sep=6pt},
}
 puis \begin{document}.
\documentclass[tikz,border=6pt]{standalone}
\usepackage[T1]{fontenc}
\usepackage[utf8]{inputenc}
\usepackage[scaled=0.95]{helvet}
\renewcommand{\familydefault}{\sfdefault}
\usepackage{amsmath,amssymb}
\usepackage[european]{circuitikz}
\usetikzlibrary{arrows.meta,positioning,calc,angles,quotes,decorations.pathreplacing,backgrounds}

% Couleurs du logo Liouaï Digital Education
\definecolor{navy}{HTML}{243743}
\definecolor{vert}{HTML}{28B78D}
\definecolor{vertfonce}{HTML}{167A5C}
\definecolor{vertclair}{HTML}{EAF7F2}
\definecolor{gris}{HTML}{A1A9AE}
\definecolor{gristexte}{HTML}{5C6870}
\definecolor{fond}{HTML}{F3F5F6}
\definecolor{jaune}{HTML}{E5C800}
\definecolor{jauneclair}{HTML}{FFFBD6}

\tikzset{
  every picture/.style={line cap=round, line join=round, color=navy, text=navy},
  fleche/.style={-{Stealth[length=2.6mm,width=2mm]}, line width=1.1pt},
  vecteur/.style={fleche, line width=1.4pt},
  titre/.style={font=\bfseries\small, text=navy},
  note/.style={font=\footnotesize, text=gristexte},
  boite/.style={draw=navy, line width=0.9pt, rounded corners=5pt, fill=fond, align=center, inner sep=6pt},
}
 puis \begin{document}.
\documentclass[tikz,border=6pt]{standalone}
\usepackage[T1]{fontenc}
\usepackage[utf8]{inputenc}
\usepackage[scaled=0.95]{helvet}
\renewcommand{\familydefault}{\sfdefault}
\usepackage{amsmath,amssymb}
\usepackage[european]{circuitikz}
\usetikzlibrary{arrows.meta,positioning,calc,angles,quotes,decorations.pathreplacing,backgrounds}

% Couleurs du logo Liouaï Digital Education
\definecolor{navy}{HTML}{243743}
\definecolor{vert}{HTML}{28B78D}
\definecolor{vertfonce}{HTML}{167A5C}
\definecolor{vertclair}{HTML}{EAF7F2}
\definecolor{gris}{HTML}{A1A9AE}
\definecolor{gristexte}{HTML}{5C6870}
\definecolor{fond}{HTML}{F3F5F6}
\definecolor{jaune}{HTML}{E5C800}
\definecolor{jauneclair}{HTML}{FFFBD6}

\tikzset{
  every picture/.style={line cap=round, line join=round, color=navy, text=navy},
  fleche/.style={-{Stealth[length=2.6mm,width=2mm]}, line width=1.1pt},
  vecteur/.style={fleche, line width=1.4pt},
  titre/.style={font=\bfseries\small, text=navy},
  note/.style={font=\footnotesize, text=gristexte},
  boite/.style={draw=navy, line width=0.9pt, rounded corners=5pt, fill=fond, align=center, inner sep=6pt},
}

\begin{document}
\begin{tikzpicture}[scale=1.0]
% ---- det > 0
\begin{scope}
\node[titre] at (1.7,3.55) {$\det A > 0$};
\draw[gris!60, line width=0.4pt] (-0.3,-0.3) grid (3.6,3.1);
\filldraw[fill=vertclair, draw=vertfonce, line width=1pt] (0,0) -- (3,1) -- (4,3) -- (1,2) -- cycle;
\draw[vecteur, vertfonce] (0,0) -- (3,1) node[below right=-1pt] {$\begin{pmatrix} a \\ c \end{pmatrix}$};
\draw[vecteur, navy] (0,0) -- (1,2) node[above left=-1pt] {$\begin{pmatrix} b \\ d \end{pmatrix}$};
\node[font=\small] at (2,1.5) {aire $= ad - bc$};
\draw[->, >=Stealth, gristexte, line width=0.8pt] (1.1,0.37) arc[start angle=18, end angle=63, radius=1.15];
\node[note, align=center] at (1.8,-0.95) {exemple : $3 \times 2 - 1 \times 1 = 5$\\ le sens est conservé};
\end{scope}
% ---- det < 0
\begin{scope}[xshift=6cm]
\node[titre] at (1.7,3.55) {$\det A < 0$};
\draw[gris!60, line width=0.4pt] (-0.3,-0.3) grid (3.6,3.1);
\filldraw[fill=jauneclair, draw=jaune!80!navy, line width=1pt] (0,0) -- (1,2) -- (4,3) -- (3,1) -- cycle;
\draw[vecteur, vertfonce] (0,0) -- (1,2) node[above left=-1pt] {$A\mathbf{e}_1$};
\draw[vecteur, navy] (0,0) -- (3,1) node[below right=-1pt] {$A\mathbf{e}_2$};
\draw[->, >=Stealth, gristexte, line width=0.8pt] (0.52,1.03) arc[start angle=63, end angle=18, radius=1.15];
\node[note, align=center] at (1.8,-0.95) {même aire, mais la figure\\ est retournée (miroir)};
\end{scope}
% ---- det = 0
\begin{scope}[xshift=12cm]
\node[titre] at (1.7,3.55) {$\det A = 0$};
\draw[gris!60, line width=0.4pt] (-0.3,-0.3) grid (3.6,3.1);
\draw[vecteur, navy, line width=2.4pt] (0,0) -- (3,2.25) node[above] {$A\mathbf{e}_2$};
\draw[vecteur, vertfonce] (0,0) -- (2,1.5) node[below right=-1pt] {$A\mathbf{e}_1$};
\node[note, align=center] at (1.8,-0.95) {colonnes alignées : le plan\\ est aplati, aire nulle};
\end{scope}
\end{tikzpicture}
\end{document}
